\documentclass[a4paper,11pt]{article}
\usepackage[utf8]{inputenc}
\usepackage[T1]{fontenc}
\usepackage[french]{babel}
\usepackage{amsmath,amsfonts,amssymb}
\usepackage{geometry}
\usepackage{fancyhdr}
\usepackage{titlesec}
\usepackage{enumitem}
\usepackage{array}

\geometry{margin=2.5cm}

% Header and footer configuration
\pagestyle{fancy}
\fancyhf{}
\fancyhead[L]{Licence 1 Informatique}
\fancyhead[R]{CERI - Avignon Université}
\fancyfoot[C]{Mathématiques Discrètes \hfill TD2 - Résolution \hfill Page \thepage/3}

% Title formatting
\titleformat{\section}
{\Large\bfseries}
{\thesection}
{1em}
{}

\begin{document}

\begin{center}
{\Huge\bfseries TD2 - Résolution} \\[0.5cm]
{\large Hiver 2026} \\[0.3cm]
\end{center}

\vspace{0.5cm}

L'objectif de ce TD est de mettre en pratique les notions de résolution.

\section{Formes normales réduites}

À partir des formules, donnez la forme normale conjonctive et réduire cette formule. Donnez les contre-modèles de la formule.

\begin{enumerate}
\item $y \Leftrightarrow (x \wedge y)$

\item $(x \Leftrightarrow (y \vee z)) \wedge (z \Rightarrow (x \vee y)) \wedge \neg(y \vee \neg x)$

\item $(x \wedge y) \vee (\neg x \wedge y \wedge \neg z)$

\item $(\neg p \Rightarrow q) \wedge ((p \Rightarrow r) \vee (q \Rightarrow p)) \wedge (r \Rightarrow p) \wedge (\neg q \Rightarrow \neg(p \wedge q)) \wedge (r \Rightarrow (p \vee q)) \wedge ((p \Rightarrow r) \vee (r \Rightarrow p)) \wedge (\neg q \Rightarrow \neg r)$

\item $(x \vee y \vee z) \wedge (\neg x \Rightarrow y) \wedge (z \Leftrightarrow y)$
\end{enumerate}

\section{Identifier des inférences}

Pour chacun des cas, utilisez un formalisme logique pour traduire les hypothèses et la conclusion et précisez la règle d'inférence utilisée pour obtenir cette conclusion à partir des hypothèses.

\begin{enumerate}
\item \emph{``Alice étudie les mathématiques.''}\\
On en déduit qu'Alice étudie les mathématiques ou l'informatique.

\item \emph{``Bob étudie les mathématiques et l'informatique.''}\\
On en déduit que Bob étudie les mathématiques.

\item \emph{``S'il pleut alors la piscine est fermée. Aujourd'hui il pleut.''}\\
On en déduit que la piscine est fermée.

\item \emph{``S'il neige, l'université sera fermée. L'université est ouverte aujourd'hui.''}\\
On en déduit qu'il ne neige pas.

\item \emph{``Quand je me baigne, je reste trop longtemps au soleil. Si je reste au soleil trop longtemps, alors j'ai des coups de soleil.''}\\
On en déduit que si je vais me baigner, alors j'aurai des coups de soleil.
\end{enumerate}

\section{Déductions}

Pour chacun des cas, utilisez les règles d'inférence afin de détailler l'argumentation permettant d'obtenir la conclusion donnée.

\begin{enumerate}
\item Déduire $c$ à partir des hypothèses :
\begin{itemize}
\item $\neg d$
\item $(\neg q \vee b)$
\item $(b \Rightarrow d)$
\item $(a \vee c)$
\end{itemize}

\item Déduire $e$ à partir des hypothèses :
\begin{itemize}
\item $(\neg a \wedge b)$
\item $(d \Rightarrow c)$
\item $(\neg e \Rightarrow d)$
\item $(c \Rightarrow a)$
\end{itemize}

\item Déduire $a \wedge b$ à partir des hypothèses :
\begin{itemize}
\item $(\neg e \Rightarrow \neg d)$
\item $(\neg b \Rightarrow d)$
\item $(a \vee c)$
\item $(\neg a \wedge \neg c)$
\end{itemize}

\item Déduire $d$ à partir des hypothèses :
\begin{itemize}
\item $(a \Rightarrow b)$
\item $\neg c$
\item $(a \vee d)$
\item $(b \Rightarrow c)$
\end{itemize}

\item Déduire $\neg a$ à partir des hypothèses :
\begin{itemize}
\item $(d \Rightarrow b)$
\item $(c \wedge \neg b)$
\item $(a \Rightarrow d)$
\end{itemize}
\end{enumerate}

\section{Déductions à partir de formes normales conjonctives}

Pour chacune des formules, donnez la forme réduite puis utilisez les règles d'inférence afin de détailler l'argumentation permettant d'obtenir la conclusion donnée.

\begin{enumerate}
\item $(\neg a \vee b \vee c) \wedge (b \vee c \vee d \vee \neg b) \wedge (b \vee c) \wedge (a \vee \neg b) \wedge (\neg a \vee b \vee c) \wedge (\neg a \vee c \vee \neg a)$

Déduire $c$ à partir de cette formule

\item $(\neg b \vee c \vee \neg d \vee b) \wedge (a \vee b \vee c) \wedge (\neg a \vee b) \wedge (\neg b \vee c \vee \neg d) \wedge (\neg c \vee e) \wedge (a \vee b) \wedge (\neg a \vee b \vee c) \wedge (\neg a \vee b \vee c \vee a)$

Déduire $\neg d \vee e$ à partir de cette formule

\item $(a \vee c) \wedge (\neg a \vee c) \wedge (a \vee b \vee c) \wedge (b \vee \neg c) \wedge (\neg c \vee d) \wedge (a \vee b \vee \neg a) \wedge (b \vee \neg c \vee a) \wedge (\neg a \vee e \vee a)$

Déduire $b \vee d$ à partir de cette formule

\item $(a \vee b \vee \neg a \vee d) \wedge (a \vee \neg b \vee c) \wedge (a \vee b \vee \neg c) \wedge (\neg c \vee d) \wedge (\neg \neg c \vee a \vee b \vee c) \wedge (\neg c \vee d) \wedge (\neg a \vee c \vee a) \wedge (d \vee \neg d)$

Montrez que cette formule est insatisfaisable. Autrement dit, montrez qu'on peut déduire $\bot$ de cette formule.

\item $(\neg a \vee b \vee c) \wedge (a \vee \neg c) \wedge (a \vee c \vee \neg b) \wedge (\neg b \vee \neg a) \wedge (a \vee b) \wedge (\neg a \vee e \vee \neg a)$

Montrez que cette formule est insatisfaisable. Autrement dit, montrez qu'on peut déduire $\bot$ de cette formule.
\end{enumerate}

\section{Modéliser et déduire}

Pour chacun des cas, utilisez un formalisme logique pour traduire les hypothèses données, puis utilisez les règles d'inférence pour produire une argumentation afin de conclure.

\begin{enumerate}
\item \emph{``Il pleut ou il fait beau ou il neige. S'il pleut, je prends ma voiture. S'il fait beau je prends mon vélo et s'il neige je prends mes skis.''}

Montrez qu'on peut en déduire : \emph{``Je prends ma voiture, ou je prends mon vélo ou je prends mes skis.''}

\item \emph{``Si tu m'envoies un mail, je finirai d'écrire le code. Si tu ne m'envoies pas de mail, j'irai me coucher tôt. Si je vais me coucher tôt, je me réveillerai en forme.''}

Montrez qu'on peut en déduire : \emph{``Si je ne finis pas d'écrire le code, alors je me réveillerai en forme.''}

\item \emph{``Si je fais de l'athlétisme, j'ai des courbatures. Je prends un bain quand j'ai des courbatures. Je n'ai pas pris de bain.''}

Montrez qu'on peut en déduire : \emph{``Je n'ai pas fait d'athlétisme.''}

\item \emph{``Lorsque j'hallucine, je vois des éléphants roses. Je suis en train de rêver ou j'hallucine. Je bois un grand verre d'eau quand je vois des éléphants roses. Je ne rêve pas.''}

Montrez qu'on peut en déduire : \emph{``Je bois un grand verre d'eau si je ne rêve pas.''}

\item Dans une maison hantée, les esprits se manifestent sous deux formes différentes : des chants discordants ou des rires moqueurs. Il est possible d'identifier un esprit en écoutant attentivement le comportement des esprits, en jouant de l'orgue ou en brûlant de l'encens de la façon suivante :
\begin{itemize}
\item Les chants ne se font pas entendre, à moins que nous jouions de l'orgue sans que les rires ne se fassent entendre.
\item Si nous brûlons de l'encens, les rires se font entendre si et seulement si les chants se font entendre.
\end{itemize}

Actuellement, les chants se font entendre et les rires sont silencieux.

Montrez qu'on peut en déduire : \emph{``Nous jouons de l'orgue et ne brûlons pas d'encens.''}

Pour cela, vous noterez que ``$x$ à moins que $y$'' se traduit par la formulation logique $(\neg x \Leftrightarrow y)$.
\end{enumerate}

\section{Analyser des raisonnements}

Pour chacun des cas, précisez si le raisonnement est correct en justifiant votre réponse.

\begin{enumerate}
\item ``Si je réussis mes examens ou que j'ai assez d'argent, je m'offrirai une guitare. Je me suis offert une guitare.''

Peut-on en déduire que j'ai réussi mes examens ? Peut-on en déduire que j'ai assez d'argent ?

\item ``Si je réussis mes examens et que j'ai assez d'argent, je m'offrirai une guitare. J'ai assez d'argent mais je ne me suis pas offert de guitare.''

Peut-on en déduire que je n'ai pas réussi mes examens ?

\item ``Si je réussis mes examens ou que j'ai assez d'argent, je m'offrirai une guitare. J'ai réussi mes examens.''

Peut-on en déduire que je me suis offert de guitare ?

\item ``Si je réussis mes examens ou que j'ai assez d'argent, je m'offrirai une guitare. Je n'ai pas réussi mes examens et je n'ai pas assez d'argent.''

Peut-on en déduire que je ne me suis pas offert de guitare ?

\item ``Si je réussis mes examens et que j'ai assez d'argent, je m'offrirai une guitare. J'ai réussi mes examens.''

Peut-on en déduire que si j'ai assez d'argent alors je m'offrirai une guitare ?
\end{enumerate}

\section{Pour aller plus loin}

Le problème des Huit Reines est le suivant : est-il possible de placer 8 reines sur un échiquier $8 \times 8$ sans qu'aucune reine n'en menace une autre ? C'est-à-dire qu'il faut placer au maximum une reine par ligne, au maximum une reine par colonne et au maximum une reine par diagonale.

Considérons une version plus simple où on souhaite placer seulement 4 reines sur un échiquier $4 \times 4$. On utilisera les variables logiques suivantes :
\begin{itemize}
\item $x_{i,j}$ : ``Une reine se trouve dans la case $i,j$''
\end{itemize}

On notera les cases de la façon suivante :

\begin{center}
\begin{tabular}{|c|c|c|c|}
\hline
1,1 & 1,2 & 1,3 & 1,4 \\
\hline
2,1 & 2,2 & 2,3 & 2,4 \\
\hline
3,1 & 3,2 & 3,3 & 3,4 \\
\hline
4,1 & 4,2 & 4,3 & 4,4 \\
\hline
\end{tabular}
\end{center}

\begin{enumerate}
\item Formalisez le fait qu'exactement une reine doit se trouver sur la première ligne.

\item Généralisez pour toutes les lignes, toutes les colonnes, toutes les diagonales.

\item En utilisant les hypothèses établies dans les questions précédentes, montrez que si une reine est placée dans la case (1,1), alors il n'y a pas de solution.
\end{enumerate}

\end{document}
